\begin{abstract}
Automatic post-editing can improve a machine translation draft while also introducing unnecessary or unsupported changes. We study this trade-off for English--Tagalog and English--Spanish in both directions, using a sentence-split NLLB-200 3.3B baseline and two post-editors: Llama 3 8B Instruct and GPT-4o Mini. Three strategies provide a source and draft alone, add one retrieved development exemplar, or require structured error diagnostics before editing. This development manuscript reports a completed pilot on the same 100 aligned FLORES+ development IDs, yielding 2,400 system outcomes. GPT retrieved one-shot improves English-to-Tagalog COMET by 1.463 points, with a descriptive paired 95\% interval of [0.710, 2.396]. Llama reason-then-edit requires baseline fallback on 227 of 400 outcomes, compared with eight for GPT; these failures remain in the evaluation. Source-based inspection also identifies numeric corruption despite a metric gain and an apparently acceptable paraphrase with a substantial metric penalty. The pilot therefore motivates separate measurement of average quality, preservation, and output validity. It does not establish universal language-resource effects or intrinsic benefits of reasoning. Main devtest results, confirmatory comparisons, and final protocol decisions remain pending.
\end{abstract}
